\chapter{Grundlagen}
\label{ch:grundlagen}
\section{Argumentationssysteme}
\label{sec:af}

\emph{Argumentationssysteme} (engl. \emph{argumentation frameworks}, AFs)
modellieren Argumente als abstrakte Einheiten und erfassen ausschließlich die
zwischen ihnen bestehenden Angriffe. Die grundlegenden Begriffe gehen auf
\textcite{dung:1995:af} zurück. Die labellingbasierte Darstellung und die
Notation folgen \textcite{bluemel:2025:independence-aba}.

\begin{Definition}
    \label{def:af}
    Ein \emph{Argumentationssystem} ist ein Paar \(F=(A,R)\), wobei \(A\)
    eine Menge von Argumenten und \(R\subseteq A\times A\) eine
    binäre Angriffsrelation ist. Für \((a,b)\in R\) heißt \(a\) ein
    \emph{Angreifer} von \(b\). Falls die Zuordnung zu einem bestimmten
    Argumentationssystem explizit angegeben werden muss, schreiben wir
    \(A_F\) und \(R_F\).
\end{Definition}

Im Folgenden werden ausschließlich endliche Argumentationssysteme betrachtet.
\begin{Definition}
    \label{def:attack/defense}
Für \(a\in A\) und \(S\subseteq A\) bezeichnen
\[
    a^+ = \{b\in A \mid (a,b)\in R\},
\]
\[
    a^- = \{b\in A \mid (b,a)\in R\},
\]
\[
    S^+ = \{b\in A \mid \exists a\in S\colon (a,b)\in R\},
\]
\[
    S^- = \{b\in A \mid \exists a\in S\colon (b,a)\in R\}
\]
die von \(a\) beziehungsweise \(S\) angegriffenen Argumente sowie deren
Angreifer.
\end{Definition}

\begin{Definition}
    \label{def:defense}
    Sei \(F=(A,R)\) ein Argumentationssystem, \(S\subseteq A\) und
    \(a\in A\). Die Menge \(S\) \emph{verteidigt} \(a\), wenn sie jeden
    Angreifer von \(a\) angreift, das heißt, genau dann wenn \(a^-\subseteq S^+\).
\end{Definition}

Insbesondere wird jedes unangreifbare Argument durch die leere Menge
verteidigt. Eine Menge \(S\) verteidigt eine Menge \(T\subseteq A\), wenn
\(S\) jedes Argument aus \(T\) verteidigt.

Eine Argumentationssemantik legt fest, welche Bewertungen der Argumente eines
Argumentationssystems zulässig sind. In der labellingbasierten Darstellung
erhält jedes Argument genau einen von drei Akzeptanzstatus.

\begin{Definition}
    \label{def:labelling}
    Sei \(F=(A,R)\) ein Argumentationssystem und
    \(\mathcal{L}=\{\mathsf{in},\mathsf{out},\mathsf{und}\}\). Ein
    \emph{Labelling} von \(F\) ist eine Funktion
    \(\lambda\colon A\to\mathcal{L}\). Dabei steht \(\mathsf{in}\) für
    akzeptiert, \(\mathsf{out}\) für abgelehnt und \(\mathsf{und}\) für
    unentschieden. Für \(\ell\in\mathcal{L}\) schreiben wir
    \[
        \ell(\lambda)=\{a\in A \mid \lambda(a)=\ell\}.
    \]
    Somit bilden \(\mathsf{in}(\lambda)\),
    \(\mathsf{out}(\lambda)\) und \(\mathsf{und}(\lambda)\) eine Partition
    von \(A\).
\end{Definition}

\begin{Definition}
    \label{def:complete-labelling}
    Ein Labelling \(\lambda\colon A\to\mathcal{L}\) heißt
    \emph{vollständig}, wenn für jedes \(a\in A\) gilt:
    \[
        \lambda(a)=\mathsf{in}
        \quad\Longleftrightarrow\quad
        \lambda(b)=\mathsf{out}\text{ für alle }b\in a^-,
    \]
    \[
        \lambda(a)=\mathsf{out}
        \quad\Longleftrightarrow\quad
        \lambda(b)=\mathsf{in}\text{ für mindestens ein }b\in a^-.
    \]
    Ein Argument erhält folglich genau dann das Label \(\mathsf{und}\), wenn
    keiner seiner Angreifer mit \(\mathsf{in}\) gelabelt ist, aber nicht alle
    seine Angreifer mit \(\mathsf{out}\) gelabelt sind.
\end{Definition}

\begin{Definition}
    \label{def:labelling-semantics}
    Sei \(F=(A,R)\) ein Argumentationssystem. Eine \emph{Semantik}
    \(\sigma\) ordnet \(F\) eine Menge \(\Lambda_\sigma(F)\) von
    Labellings zu. Die Elemente von \(\Lambda_\sigma(F)\) heißen
    \(\sigma\)-Labellings von \(F\).

    Sei \(\lambda\colon A\to\mathcal{L}\) ein Labelling. Dann gilt:
    \begin{itemize}
        \item \(\lambda\) ist ein \emph{Complete-Labelling}, gdw.
        \(\lambda\) \emph{vollständig} ist. Wir schreiben
       \(\lambda\in\Lambda_{\mathsf{co}}(F)\).

        \item \(\lambda\) ist ein \emph{Grounded-Labelling}, gdw.
        \(\lambda\) vollständig ist und \(\mathsf{und}(\lambda)\)
        bezüglich \(\subseteq\) maximal unter den
        \(\mathsf{und}\)-Mengen aller vollständigen Labellings von \(F\)
        ist. Wir schreiben 
        \(\lambda\in\Lambda_{\mathsf{gr}}(F)\).

        \item \(\lambda\) ist ein \emph{Preferred-Labelling}, gdw.
        \(\lambda\) vollständig ist und \(\mathsf{in}(\lambda)\)
        bezüglich \(\subseteq\) maximal unter den
        \(\mathsf{in}\)-Mengen aller vollständigen Labellings von \(F\)
        ist. Wir schreiben
        \(\lambda\in\Lambda_{\mathsf{pr}}(F)\).

        \item \(\lambda\) ist ein \emph{Stable-Labelling}, gdw.
        \(\lambda\) vollständig ist und
        \(\mathsf{und}(\lambda)=\emptyset\) gilt.
        Wir schreiben
        \(\lambda\in\Lambda_{\mathsf{st}}(F)\).
    \end{itemize}
\end{Definition}

Jedes Argumentationssystem besitzt genau ein Grounded-Labelling und
mindestens ein Complete- sowie ein Preferred-Labelling. Die Menge
\(\Lambda_{\mathsf{st}}(F)\) kann hingegen leer sein.

Für die spätere Betrachtung bedingter Unabhängigkeit werden Labellings auf
Teilmengen der Argumentmenge eingeschränkt.

\begin{Definition}
    \label{def:labelling-restriction}
    Sei \(\lambda\colon A\to\mathcal{L}\) ein Labelling und \(X\subseteq A\).
    Die \emph{Restriktion} von \(\lambda\) auf \(X\) ist das partielle
    Labelling
    \[
        \lambda|_X\colon X\to\mathcal{L},
        \qquad
        (\lambda|_X)(x)=\lambda(x)
        \quad\text{für alle }x\in X.
    \]
    Insbesondere gilt \(\lambda_1|_X=\lambda_2|_X\) gdw.
    \(\lambda_1(x)=\lambda_2(x)\) für alle \(x\in X\) gilt.
\end{Definition}

\section{Bedingte Unabhängigkeit}
\label{sec:conditional-independence}

Die folgende Darstellung der bedingten Unabhängigkeit orientiert
sich an \textcite{pearl:paz:1986:graphoids}. Die Notation folgt
\textcite{bluemel:2025:independence-aba}.

\begin{Definition}
    \label{def:dependency-model}
    Sei \(V\) eine Menge. Ein \emph{Abhängigkeitsmodell} \(I\) über \(V\)
    ist eine ternäre Relation über paarweise disjunkten Teilmengen von \(V\),
    das heißt,
    \[
        I \subseteq
        \bigl\{(X,Y,Z) \bigm| X,Y,Z\subseteq V,\,
        X\cap Y=X\cap Z=Y\cap Z=\emptyset\bigr\}.
    \]
    Für paarweise disjunkte Mengen \(X,Y,Z\subseteq V\) schreiben wir
    \[
        X\perp_I Y\mid Z \text{ gdw. }
        (X,Y,Z)\in I
    \]
    und sagen, dass \(X\) und \(Y\) bezüglich \(I\) \emph{bedingt
    unabhängig gegeben \(Z\)} sind. Die Menge \(Z\) heißt
    \emph{Bedingungsmenge}, und \(X\perp_I Y\mid Z\) heißt eine
    \emph{Unabhängigkeitsaussage}. Gilt \((X,Y,Z)\notin I\), so heißen
    \(X\) und \(Y\) \emph{bedingt abhängig gegeben \(Z\)}. Wir
    schreiben \(X\not\perp_I Y\mid Z\).

    Ist das Abhängigkeitsmodell aus dem Kontext eindeutig,
    wird der Index \(I\) weggelassen. Für Elemente \(x,y\in V\) stehen
    \(x\perp y\mid Z\) und \(x\not\perp y\mid Z\) kurz für
    \(\{x\}\perp\{y\}\mid Z\) beziehungsweise
    \(\{x\}\not\perp\{y\}\mid Z\).
\end{Definition}

\begin{Definition}
    \label{def:semi-graphoid}
    Ein Abhängigkeitsmodell \(I\) über \(V\) heißt \emph{Semi-Graphoid},
    wenn für alle paarweise disjunkten Mengen
    \(X,Y,W,Z\subseteq V\) die folgenden
    \emph{Semi-Graphoid-Axiome} erfüllt sind:
    \begin{align*}
        \text{Symmetrie:}\qquad
        &X\perp Y\mid Z
        &&\Longrightarrow Y\perp X\mid Z,\\
        \text{Dekomposition:}\qquad
        &X\perp (Y\cup W)\mid Z
        &&\Longrightarrow X\perp Y\mid Z,\\
        \text{Schwache Vereinigung:}\qquad
        &X\perp (Y\cup W)\mid Z
        &&\Longrightarrow X\perp Y\mid (Z\cup W),\\
        \text{Kontraktion:}\qquad
        &\left.
          \begin{aligned}
              X&\perp Y\mid Z,\\[-0.2em]
              X&\perp W\mid (Z\cup Y)
          \end{aligned}
          \right\}
        &&\Longrightarrow X\perp (Y\cup W)\mid Z.
    \end{align*}
\end{Definition}

\begin{Definition}
    \label{def:sigma-independence-af}
    Sei \(F =(A,R)\) ein Argumentationssystem, sei \(\sigma\) eine Semantik
    und seien \(X,Y,Z\subseteq A\) paarweise disjunkt. Die Menge \(X\) ist
    \emph{\(\sigma\)-unabhängig von \(Y\) gegeben \(Z\) in \(F\)},
    geschrieben
    \[
        X\perp_\sigma Y\mid Z,
    \]
    genau dann, wenn für alle
    \(\lambda_1,\lambda_2\in\Lambda_\sigma(F)\) aus
    \(\lambda_1|_Z=\lambda_2|_Z\) die Existenz eines
    \(\lambda_3\in\Lambda_\sigma(F)\) folgt, für das
    \(
        \lambda_3|_X=\lambda_1|_X,
        \lambda_3|_Y=\lambda_2|_Y,
        \lambda_3|_Z=\lambda_1|_Z=\lambda_2|_Z
    \)
    gilt. Andernfalls sind \(X\) und \(Y\) \emph{\(\sigma\)-abhängig
    gegeben \(Z\) in \(F\)} und wir schreiben
    \(X\not\perp_\sigma Y\mid Z\).
\end{Definition}

\section{Quantifizierte Boolesche Formeln}
\label{sec:qbf}

Dieser Abschnitt führt die für den weiteren Verlauf benötigten Grundlagen
quantifizierter boolescher Formeln ein. Die Darstellung orientiert sich an
\textcite{egly:woltran:2006:reasoning-qbf}.

\begin{Definition}
    \label{def:propositional-formula}
    Sei \(\mathcal{V}\) eine abzählbare Menge \emph{aussagenlogischer
    Variablen}. Die Menge der \emph{aussagenlogischen Formeln} über
    \(\mathcal{V}\) ist induktiv wie folgt definiert:
    \begin{itemize}
        \item Jede Variable \(p\in\mathcal{V}\) sowie die
        \emph{Wahrheitskonstanten} \(\top\) und \(\bot\) sind
        aussagenlogische Formeln.

        \item Sind \(\varphi\) und \(\psi\) aussagenlogische Formeln, so sind
        auch \(\neg\varphi\), \((\varphi\land\psi)\) und
        \((\varphi\lor\psi)\) aussagenlogische Formeln.
    \end{itemize}
    Die weiteren Verknüpfungen werden durch
    \[
        \varphi\to\psi := \neg\varphi\lor\psi
        \qquad\text{und}\qquad
        \varphi\leftrightarrow\psi
        := (\varphi\to\psi)\land(\psi\to\varphi)
    \]
    abgekürzt. Mit \(\operatorname{Var}(\varphi)\) bezeichnen wir die Menge
    der in \(\varphi\) vorkommenden Variablen.

    Ein \emph{Literal} ist eine Variable \(p\) oder deren Negation
    \(\neg p\). Eine \emph{Klausel} ist eine endliche Disjunktion von
    Literalen. Eine Formel liegt in \emph{konjunktiver Normalform} (KNF)
    vor, wenn sie eine endliche Konjunktion von Klauseln ist. Dabei gelten
    die Konventionen \(\bigwedge\emptyset=\top\) und
    \(\bigvee\emptyset=\bot\).
\end{Definition}

\begin{Definition}
    \label{def:boolean-assignment}
    Sei \(V\subseteq\mathcal{V}\) endlich. Eine \emph{Belegung} von \(V\)
    ist eine Funktion \(\nu\colon V\to\{0,1\}\), wobei \(1\) für
    \emph{wahr} und \(0\) für \emph{falsch} steht. Für
    \(p\in\mathcal{V}\) und \(b\in\{0,1\}\) bezeichnet
    \(\nu[p\mapsto b]\colon V\cup\{p\}\to\{0,1\}\) die Belegung mit
    \[
        \nu[p\mapsto b](q)=
        \begin{cases}
            b,      & \text{falls }q=p,\\
            \nu(q), & \text{falls }q\in V\setminus\{p\}.
        \end{cases}
    \]
    Für \(W\subseteq V\) bezeichnet \(\nu|_W\) die Einschränkung von
    \(\nu\) auf \(W\).

    Für eine aussagenlogische Formel \(\varphi\) mit
    \(\operatorname{Var}(\varphi)\subseteq V\) ist die
    \emph{Erfüllungsrelation} \(\models\) durch die folgenden Bedingungen
    definiert:
    \begin{align*}
        &\nu\models\top,
        &&\nu\not\models\bot,\\
        &\nu\models p
        &&\Longleftrightarrow \nu(p)=1,\\
        &\nu\models\neg\varphi
        &&\Longleftrightarrow \nu\not\models\varphi,\\
        &\nu\models\varphi\land\psi
        &&\Longleftrightarrow
          \nu\models\varphi\text{ und }\nu\models\psi,\\
        &\nu\models\varphi\lor\psi
        &&\Longleftrightarrow
          \nu\models\varphi\text{ oder }\nu\models\psi.
    \end{align*}
\end{Definition}

\begin{Definition}
    \label{def:qbf-syntax}
    Eine \emph{quantifizierte boolesche Formel} (QBF) über
    \(\mathcal{V}\) wird aus aussagenlogischen Formeln zusätzlich durch die
    folgende Bildungsregel erzeugt: Ist \(\Phi\) eine QBF und
    \(p\in\mathcal{V}\), so sind auch \(\forall p\,\Phi\) und
    \(\exists p\,\Phi\) QBFs. Dabei heißen \(\forall\) und \(\exists\)
    \emph{Allquantor} beziehungsweise \emph{Existenzquantor}; \(\Phi\)
    heißt ihr jeweiliger \emph{Geltungsbereich}.

    Ein Vorkommen einer Variablen \(p\) ist \emph{gebunden}, wenn es im
    Geltungsbereich eines Quantors \(\forall p\) oder \(\exists p\) liegt.
    Andernfalls ist dieses Vorkommen \emph{frei}. Mit
    \(\operatorname{FV}(\Phi)\) und \(\operatorname{BV}(\Phi)\) bezeichnen
    wir die Mengen der in \(\Phi\) frei beziehungsweise gebunden
    vorkommenden Variablen.

    Eine QBF \(\Phi\) heißt \emph{geschlossen}, wenn
    \(\operatorname{FV}(\Phi)=\emptyset\) gilt, und andernfalls
    \emph{offen}. Sie heißt \emph{quantorenfrei}, wenn sie keinen Quantor
    enthält.
\end{Definition}

\begin{Definition}
    \label{def:qbf-semantics}
    Sei \(\Phi\) eine QBF und sei
    \(\nu\colon\operatorname{FV}(\Phi)\to\{0,1\}\) eine Belegung ihrer
    freien Variablen. Die Erfüllungsrelation für QBFs erweitert
    Definition~\ref{def:boolean-assignment} um
    \begin{align*}
        \nu\models\forall p\,\Psi
        &\quad\Longleftrightarrow\quad
        \nu[p\mapsto 0]\models\Psi
        \text{ und }
        \nu[p\mapsto 1]\models\Psi,\\
        \nu\models\exists p\,\Psi
        &\quad\Longleftrightarrow\quad
        \nu[p\mapsto 0]\models\Psi
        \text{ oder }
        \nu[p\mapsto 1]\models\Psi.
    \end{align*}

    Gilt \(\nu\models\Phi\), so heißt \(\nu\) ein \emph{Modell} von
    \(\Phi\). Die QBF \(\Phi\) heißt \emph{erfüllbar}, wenn sie mindestens
    ein Modell besitzt, und \emph{unerfüllbar}, wenn sie kein Modell
    besitzt. Sie heißt \emph{gültig}, wenn
    \(\nu\models\Phi\) für jede Belegung ihrer freien Variablen gilt. Zwei
    QBFs \(\Phi\) und \(\Psi\) heißen \emph{logisch äquivalent}, wenn
    für jede Belegung
    \(\nu\colon\operatorname{FV}(\Phi)\cup
    \operatorname{FV}(\Psi)\to\{0,1\}\) gilt:
    \[
        \nu|_{\operatorname{FV}(\Phi)}\models\Phi
        \quad\Longleftrightarrow\quad
        \nu|_{\operatorname{FV}(\Psi)}\models\Psi.
    \]

    Eine geschlossene QBF \(\Phi\) heißt \emph{wahr}, wenn
    \(\emptyset\models\Phi\) gilt, und andernfalls \emph{falsch}. Somit ist
    eine geschlossene QBF genau dann erfüllbar und gültig, wenn sie wahr ist.
\end{Definition}

\begin{Definition}
    \label{def:set-quantification}
    Sei \(\Phi\) eine QBF, sei
    \(X=\{x_1,\dots,x_n\}\subseteq\operatorname{FV}(\Phi)\) endlich und
    sei \(Q\in\{\forall,\exists\}\).
    Die \emph{Quantifizierung über die Variablenmenge \(X\)} ist die
    Kurzschreibweise
    \[
        QX\,\Phi := Qx_1\cdots Qx_n\,\Phi.
    \]
    Für \(X=\emptyset\) gilt \(QX\,\Phi:=\Phi\). Da Quantoren desselben Typs
    miteinander vertauschbar sind, ist die gewählte Reihenfolge der Variablen
    aus \(X\) für den Wahrheitswert unerheblich.

    Seien \(X_1,\dots,X_m\) endliche, paarweise disjunkte Teilmengen von
    \(\operatorname{FV}(\Phi)\). Ferner verwenden wir die Kurzschreibweisen
    \[
        \forall X_1,\dots,X_m\,\Phi
        :=\forall X_1\cdots\forall X_m\,\Phi
    \]
    und
    \[
        \exists X_1,\dots,X_m\,\Phi
        :=\exists X_1\cdots\exists X_m\,\Phi.
    \]
\end{Definition}

\begin{Definition}
    \label{def:qbf-prenex-form}
    Eine QBF befindet sich in \emph{Pränexnormalform} (PNF), wenn sie die
    Form
    \[
        Q_1X_1\,Q_2X_2\cdots Q_mX_m\,\varphi
    \]
    besitzt, wobei \(m\geq 1\), die Mengen
    \(X_1,\dots,X_m\) nichtleer und paarweise disjunkt sind,
    \(Q_i\in\{\forall,\exists\}\), \(Q_i\neq Q_{i+1}\) für
    \(1\leq i<m\) gilt und \(\varphi\) quantorenfrei ist. Der Ausdruck
    \(Q_1X_1\cdots Q_mX_m\) heißt \emph{Quantorenpräfix}, und
    \(\varphi\) heißt \emph{Matrix}. Jeder Ausdruck \(Q_iX_i\) heißt
    \emph{Quantorenblock}. Eine solche QBF besitzt folglich
    \(m-1\) Quantorenwechsel.

    Liegt die Matrix \(\varphi\) zusätzlich in konjunktiver Normalform vor,
    so befindet sich die QBF in \emph{pränexer konjunktiver Normalform}
    (PCNF).
\end{Definition}

\begin{Definition}
    \label{def:bounded-qbf-problems}
    Sei \(k\geq 1\). Mit \(\mathrm{QBF}^{\forall}_k\) bezeichnen wir das
    \emph{Wahrheitsproblem} für geschlossene QBFs in Pränexnormalform mit
    genau \(k\) alternierenden Quantorenblöcken, deren erster Block
    universell quantifiziert ist. Analog bezeichnet
    \(\mathrm{QBF}^{\exists}_k\) das Wahrheitsproblem für solche QBFs,
    deren erster Block existenziell quantifiziert ist. Gefragt ist jeweils,
    ob die gegebene QBF wahr ist.
\end{Definition}

\section{Komplexitätstheoretische Grundlagen}
\label{sec:complexity}

Die folgenden Grundlagen orientieren sich an
\textcite[Abschnitt~2.2]{dvorak:2012:computational-aspects}. Die Darstellung
beschr\"ankt sich auf die im weiteren Verlauf ben\"otigten Begriffe.

\begin{Definition}
    \label{def:problem-and-decision-problem}
    Ein \emph{Problem} \(\mathcal{P}\) wird durch eine Menge
    \(\mathcal{I}_{\mathcal{P}}\) von \emph{Instanzen} und eine f\"ur jede
    Instanz festgelegte Fragestellung beziehungsweise geforderte Ausgabe
    beschrieben. Besitzt die Fragestellung ausschlie\ss lich die m\"oglichen
    Antworten \emph{ja} und \emph{nein}, so hei\ss t \(\mathcal{P}\) ein
    \emph{Entscheidungsproblem}.

    Mit \(\{0,1\}^*\) bezeichnen wir die Menge aller endlichen W\"orter
    \"uber dem bin\"aren Alphabet \(\{0,1\}\). Eine
    \emph{bin\"are Kodierung} der Instanzen ist eine injektive Abbildung
    \[
        \operatorname{enc}_{\mathcal{P}}\colon
        \mathcal{I}_{\mathcal{P}}\to\{0,1\}^*.
    \]
    Sei \(\mathcal{Y}_{\mathcal{P}}\subseteq
    \mathcal{I}_{\mathcal{P}}\) die Menge der Instanzen, f\"ur welche die
    Antwort \emph{ja} lautet. Das zugeh\"orige
    \emph{Entscheidungsproblem als Sprache} ist
    \[
        L_{\mathcal{P}}
        =\{\operatorname{enc}_{\mathcal{P}}(x)
          \mid x\in\mathcal{Y}_{\mathcal{P}}\}
        \subseteq\{0,1\}^*.
    \]
    Im Folgenden werden ein Entscheidungsproblem und die es repr\"asentierende
    Sprache miteinander identifiziert. F\"ur eine kodierte Instanz
    \(x\in\{0,1\}^*\) hei\ss t \(x\in L_{\mathcal{P}}\) eine
    \emph{Ja-Instanz} und \(x\notin L_{\mathcal{P}}\) eine
    \emph{Nein-Instanz}. Die \emph{Eingabegr\"o\ss e} ist die Wortl\"ange
    \(|x|\). Das \emph{Komplementproblem} \(\overline{\mathcal{P}}\) wird
    durch die Sprache
    \[
        \overline{L_{\mathcal{P}}}
        =\{0,1\}^*\setminus L_{\mathcal{P}}
    \]
    repr\"asentiert.
\end{Definition}

\begin{Definition}
    \label{def:turing-machine}
    Sei \(k\geq 1\). Eine \emph{deterministische
    \(k\)-Band-Turingmaschine} ist ein Tupel
    \[
        M=(Q,\Gamma,\delta,q_0,q_{\mathsf{acc}},q_{\mathsf{rej}}),
    \]
    wobei
    \begin{itemize}
        \item \(Q\) eine endliche Menge von Zust\"anden ist,
        \item \(\Gamma\) ein endliches \emph{Bandalphabet} mit
        \(\{0,1,\square\}\subseteq\Gamma\) und dem
        \emph{Blanksymbol} \(\square\) ist,
        \item \(q_0\in Q\) der \emph{Startzustand} ist,
        \item \(q_{\mathsf{acc}},q_{\mathsf{rej}}\in Q\) verschiedene
        \emph{Haltezust\"ande} sind und
        \item
        \[
            \delta\colon
            \bigl(Q\setminus
            \{q_{\mathsf{acc}},q_{\mathsf{rej}}\}\bigr)\times\Gamma^k
            \to
            Q\times\Gamma^k\times
            \{\mathsf{L},\mathsf{R},\mathsf{N}\}^k
        \]
        die \emph{\"Ubergangsfunktion} ist. Dabei stehen
        \(\mathsf{L}\), \(\mathsf{R}\) und \(\mathsf{N}\) f\"ur eine
        Bewegung des jeweiligen Schreib-Lese-Kopfes nach links, nach rechts
        beziehungsweise f\"ur keine Bewegung.
    \end{itemize}

    Eine \emph{Konfiguration} von \(M\) besteht aus dem aktuellen Zustand,
    den Inhalten der \(k\) B\"ander und den Positionen ihrer
    Schreib-Lese-K\"opfe. Bei Eingabe \(x\in\{0,1\}^*\) enth\"alt das erste
    Band anfangs \(x\), alle \"ubrigen Bandzellen enthalten \(\square\), und
    die Maschine befindet sich in \(q_0\). Eine \emph{Berechnung} ist die
    durch wiederholte Anwendung von \(\delta\) erzeugte Folge von
    Konfigurationen. Sie ist \emph{akzeptierend} beziehungsweise
    \emph{verwerfend}, wenn sie in \(q_{\mathsf{acc}}\) beziehungsweise
    \(q_{\mathsf{rej}}\) endet.

    Die Maschine \(M\) \emph{entscheidet} ein Entscheidungsproblem
    \(\mathcal{P}\), wenn sie auf jeder Eingabe h\"alt und eine Eingabe
    \(x\) genau dann akzeptiert, wenn \(x\in L_{\mathcal{P}}\) gilt.

    Eine \emph{nichtdeterministische \(k\)-Band-Turingmaschine} entsteht,
    indem die \"Ubergangsfunktion durch eine \emph{\"Ubergangsrelation}
    \[
        \Delta\subseteq
        \left(
            \bigl(Q\setminus
            \{q_{\mathsf{acc}},q_{\mathsf{rej}}\}\bigr)\times\Gamma^k
        \right)
        \times
        \left(
            Q\times\Gamma^k\times
            \{\mathsf{L},\mathsf{R},\mathsf{N}\}^k
        \right)
    \]
    ersetzt wird. Ein \emph{Rechenzweig} ist eine durch jeweils eine der
    m\"oglichen Transitionen bestimmte Berechnung. Die Maschine akzeptiert
    eine Eingabe, wenn mindestens ein Rechenzweig akzeptierend ist, und sie
    verwirft die Eingabe, wenn alle Rechenzweige verwerfend sind. Sie
    entscheidet \(\mathcal{P}\), wenn alle Rechenzweige halten und sie eine
    Eingabe \(x\) genau dann akzeptiert, wenn
    \(x\in L_{\mathcal{P}}\) gilt.
\end{Definition}

\begin{Definition}
    \label{def:polynomial-running-time}
    Die \emph{Laufzeit} \(t_M(x)\) einer haltenden Turingmaschine \(M\) auf
    einer Eingabe \(x\) ist im deterministischen Fall die Anzahl ihrer
    Berechnungsschritte und im nichtdeterministischen Fall die maximale
    Anzahl von Schritten eines Rechenzweiges. Die Maschine arbeitet in
    \emph{Polynomialzeit}, wenn ein Polynom \(p\) existiert, sodass
    \[
        t_M(x)\leq p(|x|)
    \]
    f\"ur jede Eingabe \(x\in\{0,1\}^*\) gilt.
\end{Definition}

\begin{Definition}
    \label{def:oracle-machine}
    Sei \(O\subseteq\{0,1\}^*\) eine Sprache. Eine
    \emph{Turingmaschine mit Orakel \(O\)} beziehungsweise
    \emph{\(O\)-Orakelmaschine} ist eine deterministische oder
    nichtdeterministische Turingmaschine mit einem ausgezeichneten
    \emph{Anfrageband} und drei ausgezeichneten Zust\"anden
    \(q_{?}\), \(q_{\mathsf{yes}}\) und \(q_{\mathsf{no}}\). Befindet sich
    die Maschine in \(q_{?}\), so wird das auf dem Anfrageband stehende Wort
    \(y\) in einem Schritt durch das \emph{Orakel} gepr\"uft. Der n\"achste
    Zustand ist \(q_{\mathsf{yes}}\), falls \(y\in O\) gilt, und andernfalls
    \(q_{\mathsf{no}}\). Eine solche Anfrage hei\ss t
    \emph{Orakelaufruf} und wird unabh\"angig von der Komplexit\"at von
    \(O\) als ein Berechnungsschritt gez\"ahlt.
\end{Definition}

\begin{Definition}
    \label{def:p-np-conp}
    Eine \emph{Komplexit\"ats\-klasse} ist eine Menge von
    Entscheidungs\-problemen. Die folgenden Klassen sind durch
    Polynomialzeit beschr\"ankt:
    \begin{itemize}
        \item \(\mathrm{P}\) ist die Klasse aller Entscheidungsprobleme,
        die von einer deterministischen Turingmaschine in Polynomialzeit
        entschieden werden.

        \item \(\mathrm{NP}\) ist die Klasse aller Entscheidungsprobleme,
        die von einer nichtdeterministischen Turingmaschine in
        Polynomialzeit entschieden werden.
    \end{itemize}

    F\"ur eine Komplexit\"atsklasse \(\mathcal{C}\) ist ihre
    \emph{Komplementklasse}
    \[
        \operatorname{co}\mathcal{C}
        =\{\overline{\mathcal{P}}\mid \mathcal{P}\in\mathcal{C}\}.
    \]
    Insbesondere ist
    \[
        \mathrm{coNP}:=\operatorname{co}(\mathrm{NP}).
    \]
    Ein Entscheidungsproblem geh\"ort somit genau dann zu
    \(\mathrm{coNP}\), wenn sein Komplement zu \(\mathrm{NP}\) geh\"ort.
\end{Definition}

\begin{Definition}
    \label{def:polynomial-hierarchy}
    Sei \(O\subseteq\{0,1\}^*\) eine Sprache. Mit \(\mathrm{P}^O\)
    beziehungsweise \(\mathrm{NP}^O\) bezeichnen wir die Klassen der von
    einer deterministischen beziehungsweise nichtdeterministischen
    \(O\)-Orakelmaschine in Polynomialzeit entscheidbaren Probleme. F\"ur
    eine Komplexit\"atsklasse \(\mathcal{C}\) setzen wir
    \[
        \mathrm{P}^{\mathcal{C}}
        =\bigcup_{O\in\mathcal{C}}\mathrm{P}^O,
        \qquad
        \mathrm{NP}^{\mathcal{C}}
        =\bigcup_{O\in\mathcal{C}}\mathrm{NP}^O,
    \]
    und
    \[
        \mathrm{coNP}^{\mathcal{C}}
        :=\operatorname{co}\!\left(\mathrm{NP}^{\mathcal{C}}\right).
    \]

    Die Klassen der \emph{polynomiellen Hierarchie} werden f\"ur
    \(i\geq 0\) rekursiv durch
    \[
        \Sigma^{\mathrm{P}}_0
        =\Pi^{\mathrm{P}}_0
        =\mathrm{P},
        \qquad
        \Sigma^{\mathrm{P}}_{i+1}
        =\mathrm{NP}^{\Sigma^{\mathrm{P}}_i},
        \qquad
        \Pi^{\mathrm{P}}_{i+1}
        =\operatorname{co}\!\left(\Sigma^{\mathrm{P}}_{i+1}\right)
        =\mathrm{coNP}^{\Sigma^{\mathrm{P}}_i}
    \]
    definiert. Ihre Gesamtheit ist die Klasse
    \[
        \mathrm{PH}
        :=\bigcup_{i\geq 0}\Sigma^{\mathrm{P}}_i.
    \]
    F\"ur die ersten drei Ebenen gelten insbesondere
    \[
        \Sigma^{\mathrm{P}}_1=\mathrm{NP},
        \qquad
        \Pi^{\mathrm{P}}_1=\mathrm{coNP},
    \]
    \[
        \Sigma^{\mathrm{P}}_2
        =\mathrm{NP}^{\mathrm{NP}}
        =\operatorname{co}\!\left(\Pi^{\mathrm{P}}_2\right),
        \qquad
        \Pi^{\mathrm{P}}_2
        =\mathrm{coNP}^{\mathrm{NP}},
    \]
    sowie
    \[
        \Sigma^{\mathrm{P}}_3
        =\mathrm{NP}^{\Sigma^{\mathrm{P}}_2},
        \qquad
        \Pi^{\mathrm{P}}_3
        =\mathrm{coNP}^{\Sigma^{\mathrm{P}}_2}.
    \]
    Die Schreibweise
    \(\operatorname{co}\!\left(\Pi^{\mathrm{P}}_2\right)\) bezeichnet
    folglich dieselbe Klasse wie \(\Sigma^{\mathrm{P}}_2\) und keine
    zus\"atzliche Ebene der Hierarchie.
\end{Definition}

\begin{Definition}
    \label{def:polynomial-many-one-reduction}
    Seien \(\mathcal{A}\) und \(\mathcal{B}\) Entscheidungsprobleme. Eine
    \emph{polynomielle Many-one-Reduktion} beziehungsweise
    \emph{P-Reduktion} von \(\mathcal{A}\) auf \(\mathcal{B}\) ist eine
    totale Funktion
    \[
        f\colon\{0,1\}^*\to\{0,1\}^*,
    \]
    die von einer deterministischen Turingmaschine in Polynomialzeit
    berechnet wird und f\"ur jede Eingabe \(x\in\{0,1\}^*\) die Bedingung
    \[
        x\in L_{\mathcal{A}}
        \quad\Longleftrightarrow\quad
        f(x)\in L_{\mathcal{B}}
    \]
    erf\"ullt. In diesem Fall schreiben wir
    \(\mathcal{A}\leq^{\mathrm{P}}_{\mathrm{m}}\mathcal{B}\).
    Sofern nichts anderes angegeben ist, bezeichnet der Begriff
    \emph{Reduktion} im Folgenden stets eine polynomielle
    Many-one-Reduktion.
\end{Definition}

Eine solche Reduktion transformiert Ja-Instanzen in Ja-Instanzen und
Nein-Instanzen in Nein-Instanzen. Ein Verfahren f\"ur \(\mathcal{B}\) kann
daher nach der Transformation auch zur Entscheidung von \(\mathcal{A}\)
verwendet werden. Polynomielle Many-one-Reduktionen sind transitiv.

\begin{Definition}
    \label{def:hardness-and-completeness}
    Sei \(\mathcal{C}\) eine Komplexit\"ats\-klasse und sei
    \(\mathcal{B}\) ein Entscheidungs\-problem. Das Problem
    \(\mathcal{B}\) hei\ss t \emph{\(\mathcal{C}\)-hart}, wenn f\"ur jedes
    Problem \(\mathcal{A}\in\mathcal{C}\) eine Reduktion
    \[
        \mathcal{A}\leq^{\mathrm{P}}_{\mathrm{m}}\mathcal{B}
    \]
    existiert. Gilt zus\"atzlich \(\mathcal{B}\in\mathcal{C}\), so hei\ss t
    \(\mathcal{B}\) \emph{\(\mathcal{C}\)-vollst\"andig}. Der Nachweis
    \(\mathcal{B}\in\mathcal{C}\) hei\ss t
    \emph{Zugeh\"origkeitsnachweis}; der Nachweis der
    \(\mathcal{C}\)-H\"arte hei\ss t \emph{H\"artenachweis}.
\end{Definition}

Aufgrund der Transitivit\"at der Reduktion gen\"ugt es f\"ur einen
H\"artenachweis, ein bereits als \(\mathcal{C}\)-hart bekanntes Problem auf
das betrachtete Problem zu reduzieren.

\begin{Definition}
    \label{def:qbf-hardness-problems}
    Aufbauend auf Definition~\ref{def:bounded-qbf-problems} werden die
    folgenden einge\-schr\"ankten QBF-Entscheidungs\-probleme verwendet:
    \begin{description}
        \item[\(\mathrm{QBF}^{\forall}_{2,\mathrm{KNF}}\).]
        Eine Instanz ist eine geschlossene QBF
        \[
            \Phi=\forall X\,\exists Y\,\varphi
        \]
        in Pr\"anexnormalform, wobei \(X\) und \(Y\) disjunkte, nichtleere
        Variablenmengen sind und die Matrix \(\varphi\) eine Formel in
        konjunktiver Normalform \"uber \(X\cup Y\) ist. Gefragt ist, ob
        \(\Phi\) wahr ist.

        \item[\(\overline{\mathrm{QBF}^{\exists}_{3,\mathrm{KNF}}}\).]
        Eine Instanz ist eine geschlossene QBF
        \[
            \Psi=\exists X\,\forall Y\,\exists Z\,\varphi
        \]
        in Pr\"anexnormalform, wobei \(X\), \(Y\) und \(Z\) paarweise
        disjunkte, nichtleere Variablenmengen sind und die Matrix
        \(\varphi\) eine Formel in konjunktiver Normalform \"uber
        \(X\cup Y\cup Z\) ist. Gefragt ist, ob \(\Psi\) falsch ist. Der
        Querstrich bezeichnet somit das Komplement des entsprechenden
        QBF-Wahrheitsproblems.
    \end{description}
\end{Definition}

\begin{Theorem}[
    Meyer und Stockmeyer~\cite{meyer:stockmeyer:1972:equivalence-regex};
    Stockmeyer und Meyer~{\cite[Theorem~4.2]
    {stockmeyer:meyer:1973:word-problems}}
]
\label{thm:qbf-ph-completeness}
Für jedes \(k \geq 1\) ist \(\mathrm{QBF}^{\exists}_k\)
\(\Sigma^{\mathrm{P}}_k\)-vollständig und
\(\mathrm{QBF}^{\forall}_k\)
\(\Pi^{\mathrm{P}}_k\)-vollständig.
Die Vollständigkeit gilt bezüglich polynomialer
Many-one-Reduktionen. Insbesondere sind:
\begin{enumerate}
    \item \(\mathrm{QBF}^{\forall}_{2,\mathrm{KNF}}\)
    \(\Pi^{\mathrm{P}}_2\)-vollständig und

    \item \(\overline{\mathrm{QBF}^{\exists}_{3,\mathrm{KNF}}}\)
    \(\Pi^{\mathrm{P}}_3\)-vollständig.
\end{enumerate}
\end{Theorem}

Der erste H\"artenachweis f\"ur die alternations\-beschr\"ankten
QBF-Wahrheits\-probleme geht auf \textcite{meyer:stockmeyer:1972:equivalence-regex}
zur\"uck. Folglich kann \(\Pi^{\mathrm{P}}_2\)-H\"arte durch eine Reduktion
von \(\mathrm{QBF}^{\forall}_{2,\mathrm{KNF}}\) gezeigt werden. F\"ur
\(\Pi^{\mathrm{P}}_3\)-H\"arte kann entsprechend
\(\overline{\mathrm{QBF}^{\exists}_{3,\mathrm{KNF}}}\) verwendet werden;
bei dessen Ja-Instanzen ist die eingegebene QBF falsch.
